\section{QMTester Approach}

We design QMTester as a paired-execution workflow over \emph{program inputs}. The input is not a fixed circuit but a \emph{program subject}: a Qiskit builder together with the auditable semantics of its inputs (qubit roles, classical-register layout, parameters, ancilla protocol), a backend configuration, a shot budget, and a set of enabled program-level relation families. The output is a set of relation-level decisions plus a source-level detection decision after multiple-testing correction. Figure~\ref{fig:architecture} summarizes the workflow.

\begin{figure}[t]
\centering
\includegraphics[width=\columnwidth]{architecture_ai.png}
\caption{QMTester workflow: a program subject yields source and follow-up \emph{inputs}; the builder produces source and follow-up circuits; paired execution yields count vectors; canonicalization validates and remaps outputs; a statistical oracle flags relation violations.}
\label{fig:architecture}
\end{figure}

\subsection{Two-Layer Design}

QMTester is organized as two layers with a clean interface between them.

\noindent\textbf{Layer 1: relation admission and canonicalization.} This layer is subject-agnostic. It admits a candidate transformation only under its declared preconditions, builds the source and follow-up circuits, executes them under matched settings, and canonicalizes the resulting count vectors through a validated bijective map before any test is run (Sec.~\ref{sec:count-vectors}). A candidate whose canonicalization map is not bijective on the measured support, or that changes the measured-bit count, drops a register, or crosses a non-unitary boundary a family does not admit, is rejected and logged as a coverage loss rather than scored.

\noindent\textbf{Layer 2: program-level subject adapters and relation families.} This layer encodes the auditable builder-level semantics of each subject. A \texttt{ProgramSubject} adapter exposes how to construct the circuit from a given input and what each input position means (which qubit plays which role, which classical bit records which logical value, which parameter must respect which periodicity, which register is an ancilla). A \texttt{ProgramRelationCandidate} pairs a subject with a relation family and produces the follow-up \emph{input} together with the induced canonicalization map. Because the relation is expressed over inputs and the circuit is \emph{rebuilt} from the follow-up input, a defect that a buggy builder bakes into its output is exposed as a divergence between the source and follow-up outputs.

\subsection{Workflow}

We split the workflow into five stages (Fig.~\ref{fig:architecture}'s boxes further split stage one into input/validate and stage five into canonicalize/oracle/decision). Stage one normalizes the source subject metadata so that qubit roles, classical registers, measurement layout, and output-bit order are explicit. Stage two enumerates candidate follow-up inputs from the enabled relation families. Stage three checks each candidate against its relation-specific preconditions and constructs the induced canonicalization map before admitting a follow-up. Stage four rebuilds and executes the source and follow-up circuits under the same backend, shot budget, seed policy, and timeout policy. Stage five canonicalizes the observed count vectors and compares them with the statistical oracle.

\begin{algorithm}[t]
\caption{QMTester program-level detection pipeline}
\label{alg:qmtester}
\scriptsize
\begin{algorithmic}[1]
\Require subjects $\mathcal{S}$, program relations $\mathcal{R}$, backend $b$, shots $n$, seeds, timeout $\tau$, threshold $\alpha$
\Ensure pair logs, source decisions, rejected relation records
\ForAll{$P \in \mathcal{S}$}
  \State $Pairs \gets [\,]$; $Rejects \gets [\,]$
  \ForAll{$R \in \mathcal{R}$}
    \ForAll{$\iota' \in \Call{FollowUpInputs}{R,P}$}
      \If{$\Call{Admissible}{P,R,\iota'}$}
        \State $c \gets \Call{CanonMap}{P,R,\iota'}$
        \State $Pairs \gets Pairs \cup \{(R,P,\Call{Build}{P,\iota'},c)\}$
      \Else
        \State $Rejects \gets Rejects \cup \{(P,R,\iota',\Call{Reason}{P,R,\iota'})\}$
      \EndIf
    \EndFor
  \EndFor
  \State $Tests \gets [\,]$
  \ForAll{$(R,P,P',c) \in Pairs$}
    \State $x \gets \Call{Execute}{P,b,n,seeds,\tau}$
    \State $y \gets \Call{Execute}{P',b,n,seeds,\tau}$
    \State $(x,y) \gets \Call{Canonicalize}{x,y,c}$
    \State $(p,sparse) \gets \Call{TwoSampleTest}{x,y}$
    \If{$sparse$}
      \State $p \gets \Call{ExactOrResample}{x,y}$
    \EndIf
    \State $Tests \gets Tests \cup \{(R,p,sparse)\}$
  \EndFor
  \State $\alpha_{\mathrm{eff}} \gets \alpha/|Tests|$
  \State $detected \gets \exists(R,p,s) \in Tests: p < \alpha_{\mathrm{eff}}$
  \State \Call{WriteLogs}{$P,Pairs,Tests,Rejects$}
  \State \Call{WriteDecision}{$P,detected,\alpha_{\mathrm{eff}}$}
\EndFor
\end{algorithmic}
\end{algorithm}

We make the procedure intentionally conservative about invalid relations. A follow-up input that cannot be justified under the subject's role semantics, parameter domain, classical-control behavior, or canonicalization bijectivity is rejected before execution. This rejection reduces relation coverage, but it prevents QMTester from treating a self-inflicted semantic change as a program fault.

\subsection{Formal Model and Assumptions}

We use the formal model to mark exactly where QMTester relies on relation validity and where it relies on statistical testing. Let $P$ be a program subject with input $\iota$ that builds a circuit $\mathrm{Build}(P,\iota)$ whose measurement outcomes range over bitstrings $B(P)$. A program relation family $R$ induces candidate follow-up inputs $\iota'$; the predicate $A_R(P,\iota')$ records whether the relation's preconditions are satisfied and whether the induced map $c_R:B(\mathrm{Build}(P,\iota'))\rightarrow B(\mathrm{Build}(P,\iota))$ is bijective on the measured support. If $A_R(P,\iota')$ holds, QMTester admits the follow-up circuit $P'=\mathrm{Build}(P,\iota')$ and records $c_R$. Under a fixed execution setting $e$, repeated executions yield empirical count vectors $\hat{D}_{P,e}$ and $\hat{D}_{P',e}$ that estimate the corresponding measurement distributions defined by the circuit semantics and Born rule~\cite{nielsen2010quantum}.

We deliberately separate relation obligations from statistical decisions. When the admitted relation is valid for $P$ and the builder implements the intended semantics, then $D_{P,e}(b)=D_{P',e}(c_R^{-1}(b))$ for every observable bitstring $b$ after canonicalization. This is a conditional obligation, not an implementation-level proof; the artifact exposes admitted and rejected follow-up inputs so readers can audit $A_R$ on concrete subjects. For each admitted pair, QMTester tests $H_0:D_{P,e}=c_R(D_{P',e})$ against a distributional alternative.

\noindent\textbf{Remark 1 (source-level Bonferroni control).} For one source subject with $m$ admitted follow-up pairs, rejecting when any pair has $p_i<\alpha/m$ controls the source-level false-alarm probability at $\alpha$ when all $m$ nulls hold: each valid null has $\Pr[p_i<\alpha/m]\leq\alpha/m$, and the union bound gives $\Pr[\exists i:p_i<\alpha/m]\leq\alpha$. This is the standard Bonferroni bound; its only non-trivial content is the premise that the per-pair nulls are \emph{valid}, which is exactly what the empirical type-I check of Sec.~\ref{sec:oracle-behavior} probes (and finds conservative rather than uniform). The bound covers neither invalid relations nor unaudited raw logs.

\subsection{Fault Model and Scope}
\label{sec:fault-model}

We target faults that manifest as changes in the measured output distribution under a transformation of program inputs that should preserve observable behavior. Examples include permuted qubit roles in multi-controlled gates, wrong classical-bit indexing or endianness, inverted QFT direction, parameter values that break an intended periodicity, and ancillas that are not uncomputed. Such faults are common in quantum-program benchmarks because small builder-level mistakes can leave code type-correct while changing the program's observable behavior.

We do not claim coverage of all quantum bugs. A fault that preserves the measured distribution for the selected input and relation family will not be detected. A fault whose effect size is smaller than the available shot budget may also evade the statistical oracle. Faults that require semantic knowledge outside the program, such as an application-specific cost-function expectation, may need assertions or domain-specific relations. These boundaries define the method's scope rather than exceptions discovered after evaluation.

\subsection{Program-Level Metamorphic Relations}
\label{sec:relations}

We organize follow-up generation into five program-level relation families, each defined over program inputs and each carrying a documented soundness argument and a bijective canonicalization map.

\begin{itemize}
  \item \texttt{program\_input\_permutation}: permute the logical roles of input qubits in a way the subject's semantics declare covariant, and restore the output order by the inverse permutation. Detection does \emph{not} rest on symmetric-control invariance (which a correct builder satisfies trivially): in our \texttt{ccx\_role} subject the buggy builder hard-codes \texttt{ccx(0,1,2)} (target on wire~2) and ignores the requested role map, so permuting the input roles moves the state-prep gates but not the hard-coded \texttt{ccx}, and the canonicalized source/follow-up diverge; the role-covariant fix (\texttt{ccx(2,1,0)}) moves the gate too and stays equal, raising no false positive.
  \item \texttt{program\_classical\_remap}: remap the classical-register layout (order/endianness) under which results are recorded, then canonicalize back to the source layout. Targets measurement-order and endianness contracts.
  \item \texttt{program\_qft\_round\_trip}: compose the program's QFT stage with its declared inverse so the round trip is the identity on the logical input. Targets QFT-direction and rotation-sign defects.
  \item \texttt{program\_parameter\_periodicity}: shift a rotation parameter by a full intended period so the observable distribution is unchanged. The admission check is period-aware: a $2\pi$ shift is sound for an uncontrolled Pauli rotation (the extra $-1$ is a global phase) but a \emph{controlled} Pauli rotation has period $4\pi$ (the $-1$ becomes relative), so a $2\pi$ declaration on a controlled rotation is rejected as an invalid context rather than tested---a confusion documented in production Qiskit~\cite{qiskit13670}.
  \item \texttt{program\_ancilla\_uncompute}: require that ancilla and feedback registers return to a clean state, so adding or completing the declared uncompute leaves the logical output invariant. Targets missing/incorrect uncompute, teleportation feedback, and Grover ancilla protocols.
  \item \texttt{program\_register\_shape}: a \emph{complementary static check}, not a distribution relation---a builder asked to report in its declared $n$-bit register must emit that shape, which the prevalent \texttt{measure\_all} defect violates by doubling the output to $2n$ bits. Needing no statistics, it is the one case a deployable static analyzer also reaches (Table~\ref{tab:program-baselines}); it brings \texttt{qiskit\_github\_12} and the 37 LintQ \texttt{measure\_all} findings (\S\ref{sec:threats}) in scope.
\end{itemize}

\begin{table*}[t]
\centering
\caption{Program-level relation-validity obligations for QMTester. Each relation is defined over \emph{program inputs}; the follow-up circuit is rebuilt from the follow-up input and its outputs are remapped through a validated bijective canonicalization map before testing.}
\label{tab:relation-validity}
\scriptsize
\setlength{\tabcolsep}{2.2pt}
\begin{tabular}{@{}p{0.15\textwidth}p{0.17\textwidth}p{0.14\textwidth}p{0.15\textwidth}p{0.16\textwidth}p{0.13\textwidth}@{}}
\toprule
\textbf{Relation family} & \textbf{Preconditions} & \textbf{Invariant} & \textbf{Target bug classes} & \textbf{Invalid contexts} & \textbf{Implementation checks} \\
\midrule
\texttt{program\_input\_permutation} & Subject declares a covariant permutation of input qubit roles (e.g.\ symmetric controls); inverse permutation restores output order. & Canonicalized bitstring distribution matches the source after the inverse role permutation. & Qubit-role bookkeeping in multi-controlled and symmetric constructions. & Permutations the subject does not declare covariant; untracked classical side effects. & Validate role permutation, bijective inverse, and classical-bit remapping. \\
\texttt{program\_classical\_remap} & Classical-register layout (order/endianness) is remappable and canonically invertible to the source layout. & Logical outcome distribution is invariant after remapping classical bits to source order. & Measurement-order and endianness contracts. & Remaps that drop or duplicate classical bits, or change the measured-bit count. & Validate bijective register map, key length, and flatten/endianness handling. \\
\texttt{program\_qft\_round\_trip} & Subject exposes a QFT stage with a declared inverse over the same ordered qubits. & QFT composed with its inverse acts as the identity on the logical input. & QFT-direction and rotation-sign defects. & Approximate or truncated QFT without a matching inverse; mismatched qubit order. & Check QFT/inverse pairing, qubit order, and round-trip identity. \\
\texttt{program\_parameter\_periodicity} & Rotation parameter has a declared period $T$; shifting by $T$ preserves the observable. & Measurement distribution is invariant under a full-period parameter shift. & Angle handling and periodicity contracts (also drives hard mutants). & Parameters whose period is not declared or is broken by the subject; non-periodic gates. & Confirm declared period and reject shifts outside the periodic domain. \\
\texttt{program\_ancilla\_uncompute} & Ancilla/feedback registers are declared and must return to a clean state. & Completing the declared uncompute leaves the logical output invariant. & Missing/incorrect uncompute, teleportation feedback, Grover ancilla protocols. & Transformations that leave an ancilla entangled or cross a non-unitary boundary the family does not admit. & Verify ancilla cleanliness, feedback handling, and uncompute completion. \\
\bottomrule
\end{tabular}
\end{table*}


Our design goal is distinct builder-level contracts, not a universal relation. We distinguish families by \emph{why} circuit-level testing misses them. \texttt{program\_input\_permutation}, \texttt{program\_classical\_remap}, and largely \texttt{program\_ancilla\_uncompute} are \emph{structurally} builder-only: no circuit-level metamorphic \emph{self}-transform recovers the qubit-role mapping, register layout/endianness, or ancilla identity, because it rewrites the emitted circuit rather than rebuilding from a different input. By contrast \texttt{program\_parameter\_periodicity}, and \texttt{program\_qft\_round\_trip} where the QFT block is identifiable, are \emph{circuit-expressible}: a period shift is a local gate edit a circuit-level tool could in principle apply, so MorphQ's $0$ there reflects a \emph{catalog gap}, not structural invisibility---the necessity claim holds strictly for the structurally builder-only families, which Proposition~1 formalizes. (The one mined QPE subject is \emph{incidentally} flagged by a circuit-level swap rewrite---perturbing its wiring, not recovering the role map---so the structural argument is intact.)

\noindent\textbf{Proposition~1 (Necessity for builder-ignores-input defects).}
\label{prop:necessity}
Let a \emph{self-transform} be any MorphQ-class circuit transform---identity insertion, equivalence-preserving relabeling, phase or equivalence rewriting---that rewrites the emitted circuit $C=\mathrm{Build}(P,\iota)$ into a $C'$ preserving its unitary, hence its distribution ($D_{C'}=D_C$). For a \emph{builder-ignores-input} defect (e.g.\ \texttt{ccx\_role}), where the buggy builder hard-codes gates independently of $\iota$, no self-transform yields a source/follow-up divergence: it never re-invokes the builder, so the pair agrees by construction. Exposing the defect requires \emph{rebuilding} from a different input $\iota'$---the input-level relation $A_R$. For ancilla-uncompute the obstruction is the same, the missing uncompute crossing the non-unitary boundary (\S\ref{sec:nonunitary}) an equivalence-preserving rewrite may not cross. This is \emph{not} the false claim that information is lost on emission: a reference-oracle differential recovers all seven (Table~\ref{tab:program-baselines}) by comparing $C$ against a \emph{different} correct circuit---what no self-transform synthesizes from $C$ alone, but rebuilding from $\iota'$ supplies oracle-free. Table~\ref{tab:relation-validity} is part of the method: each row implies negative tests for the generator (e.g.\ \texttt{program\_classical\_remap} must reject non-bijective maps). A purely circuit-level catalog is retained only as a diagnostic baseline (\S\ref{sec:diagnostic}), where it shows no advantage over MorphQ.

\subsection{Statistical Oracle}
\label{sec:statistical-oracle}

After output canonicalization, we treat observed shot counts from $P$ and admitted follow-up $P'$ as samples from categorical distributions $D_P$ and $D_{P'}$ over the aligned support $B(P)$, testing $H_0: D_P = D_{P'}$ against the alternative that the distributions differ. \emph{The oracle is a Monte-Carlo permutation two-sample test}: we draw $B=10{,}000$ relabelings of the pooled shots and read the p-value off the permutation reference (\S\ref{sec:oracle-behavior}; sparse-cell mechanics below). An asymptotic $\chi^2$ fast path exists for dense supports but is not the operative branch on our exponentially sparse quantum outputs (zero cell-merges across all admitted pairs; the permutation branch decides $\geq 85.7\%$ of pairs), so it is reported, not relied on. The permutation statistic is the pooled-denominator two-sample form
\[
\begin{gathered}
T = \sum_{k=1}^{K} \frac{(x_k - e_k)^2 + (y_k - f_k)^2}{e_k + f_k},\\[2pt]
e_k=\tfrac{n_x(x_k+y_k)}{n_x+n_y},\quad f_k=\tfrac{n_y(x_k+y_k)}{n_x+n_y},
\end{gathered}
\]
for source/follow-up shot totals $n_x, n_y$. $T$ is a monotone transform of the standard Pearson two-sample statistic~\cite{pearson1900criterion}---one-half of it under matched totals $n_x=n_y$---so the permutation reference, which decides the vast majority of pairs, is unaffected by the choice. (The $\chi^2_{K'-1}$ fast path acts on the \emph{standard} Pearson statistic; per above it is untaken here.)

\noindent\textbf{Sparse-cell handling.}
Quantum outputs are sparse because the support $\{0,1\}^q$ grows exponentially with the qubit count $q$. The chi-squared approximation degrades when too many cells have small expected counts. We use a two-stage rule with explicit thresholds rather than a single heuristic:
\begin{enumerate}
  \item \emph{Cell merging.} Sort cells by $e_k$ ascending. Merge any cell with $e_k<5$ into a single ``other'' bin (Cochran's rule). If after merging the number of effective bins $K'$ satisfies $K'\geq 5$ \emph{and} at least 80\% of bins have $e_k\geq 5$, we accept the asymptotic chi-squared p-value.
  \item \emph{Exact / resampling fallback.} Otherwise we mark the pair \texttt{sparse=true} and use a Monte-Carlo permutation test on the pooled multinomial: we draw $B=10{,}000$ permutations of the $n_x+n_y$ pooled shots into two groups of size $n_x$ and $n_y$, recompute $\chi^2$, and report $\hat{p} = (1+\#\{b: \chi^2_b \geq \chi^2_{\mathrm{obs}}\})/(B+1)$. The branch is logged per pair so readers can recompute either from raw counts.
\end{enumerate}
As noted above, the fast path is not the operative branch on our corpus (zero cell-merges; $K'=K$ throughout), and the permutation branch decides $\geq 85.7\%$ of pairs (\S\ref{sec:oracle-behavior}).

\noindent\textbf{Multiple-testing correction.}
QMTester applies Bonferroni correction over all admitted follow-up pairs generated for the same source subject. With $m$ pairs and family-wise threshold $\alpha=0.05$, the default per-pair threshold is $\alpha/m$; step-down Holm correction is reported as a sensitivity check but is not mixed with the default decision rule~\cite{holm1979simple}. We use the source subject rather than the whole benchmark as the correction unit. This choice controls the chance that one subject is flagged because many follow-up variants were tried, while avoiding an overly conservative threshold that couples unrelated subjects.

\noindent\textbf{Noisy-simulator calibration.}
On noisy backends, two executions of the \emph{same} source circuit under independent seeds can produce statistically distinguishable count vectors purely from backend variability. We separate that variability from a real follow-up deviation by executing $P$ under $K=20$ independent seeds at shot budget $n$, computing the $\binom{K}{2}=190$ pairwise null statistics after the same canonicalization and sparse-cell rule, and using the $(1-\alpha/m)$-quantile as the calibrated rejection threshold $\hat\tau_P$ for that subject. The calibration runs are logged separately so readers can recompute $\hat\tau_P$ from the 20 raw count vectors. This procedure is a pragmatic complement to learned noise-mitigation methods~\cite{muqeet2024noise} and does not require a noise model.

We report all detection and false-positive rates with Wilson 95\% confidence intervals~\cite{wilson1927probable}, and the artifact records, for every admitted pair, the raw count vectors, $\chi^2$ statistic, asymptotic p-value, sparse-cell decision, fallback p-value (if applicable), and Bonferroni-corrected threshold. The prototype is a Python workflow over Qiskit's \texttt{QuantumCircuit}; all reported claims are simulator-based, and device evaluation is out of scope (Sec.~\ref{sec:threats}).
